\documentclass[10pt]{article}
\usepackage[a4paper,margin=21mm]{geometry}
\usepackage[utf8]{inputenc}
\usepackage{amsmath,amssymb,amsthm}
\usepackage{array}
\IfFileExists{booktabs.sty}{\usepackage{booktabs}}{%
  \newcommand{\toprule}{\hline}
  \newcommand{\midrule}{\hline}
  \newcommand{\bottomrule}{\hline}}
\IfFileExists{microtype.sty}{\usepackage{microtype}}{}
\IfFileExists{hyperref.sty}{%
  \IfFileExists{pdftexcmds.sty}{\usepackage[hidelinks]{hyperref}}{}}{}
\setlength{\parindent}{1.15em}
\setlength{\parskip}{0.12em}
\setlength{\tabcolsep}{5pt}
\renewcommand{\arraystretch}{1.12}

\theoremstyle{plain}
\newtheorem{theorem}{Theorem}
\newtheorem{lemma}[theorem]{Lemma}
\newtheorem{proposition}[theorem]{Proposition}
\newtheorem{corollary}[theorem]{Corollary}
\theoremstyle{definition}
\newtheorem{definition}[theorem]{Definition}
\theoremstyle{remark}
\newtheorem{remark}[theorem]{Remark}
\newcommand{\W}{\mathbb Z/2^{16}\mathbb Z}
\newcommand{\Zb}[1]{\mathbb Z/2^{#1}\mathbb Z}
\newcommand{\xor}{\mathbin{\oplus}}
\newcommand{\OR}{\mathbin{\mathsf{OR}}}
\newcommand{\E}{\mathsf E}
\newcommand{\D}{\mathsf D}
\newcommand{\rol}{\operatorname{rol}_{16}}
\newcommand{\hi}{\operatorname{hi}}
\newcommand{\midb}{\operatorname{mid}}
\newcommand{\lane}{\operatorname{lane}}
\newcommand{\topx}{\operatorname{top}_{12}}
\newcommand{\code}[1]{\texttt{#1}}

\title{Invariant Quotients in SEPAR: Practical Key Recovery with Chosen IVs and Repeated Resets}
\author{Noah Ostle\\University of Technology Sydney}
\date{13 September 2026}

\begin{document}
\maketitle

\begin{abstract}
We analyze the reference implementation of SEPAR, a stateful cipher with a
256-bit key and a 16-bit data path.  Because four assignments in its diffusion
layer are executed sequentially, the implemented map is
\(L(a,b,c,d)=(b\xor c,b,a,a\xor d)\).  Every keyed stage is consequently
triangular: one nibble evolves autonomously, while the high byte, middle byte,
and top 12 bits form exact permutation quotients.  Only
\((16!)^{4129}\) of the \(65536!\) permutations on a word have this form.
The key schedule further creates large, non-uniform fibres, the eight stages
use disjoint 32-bit key pairs, and the nominal ninth state word is determined
by the other eight.

We use these properties to factor key recovery from the outside in.  A
fixed-IV prefix experiment orders candidates for the last stage.  Thereafter
the whole attack is described by one pivot,
\(C_t=R_t+2^8h_{t+1}\): invert a candidate stage, remove the appropriate state
byte, and recurse.  We prove a bidirectional carry bound that never rejects
the true path, and finish stage one with an exact constant-difference test.
Rankings affect only traversal order.  If every score tier is visited and
early stopping is removed, the search reaches the original key for every
instance; a reachable example in which the true outer lane ranks second shows
why no rank-one assumption is permissible.

The threaded C proof of concept targets a repeated-reset, chosen-IV oracle.
With one score tier at each of its three ranking levels, it recovered 95 of
100 preselected random keys, conditional on the correct last-stage pair.  The
exact one-sided 95\% lower confidence bound on success is 0.8977.  Successful
runs took a median of 181.1 seconds on the test host.  The matching outer-stage
experiment selected the correct pair first for all 100 keys, although this is
an empirical observation rather than a universal claim.
\end{abstract}

\noindent\textbf{Keywords:} lightweight cryptography; invariant quotient;
chosen IV; key recovery; SEPAR.

\section{Introduction}

SEPAR combines eight keyed 16-bit permutations with eight state words and an
LFSR-like word \cite{VahiJassbi2020}.  Its security discussion is based on
local S-box properties, active-S-box counts, a nominal 144-bit state, and
statistical tests of the output.  Later work multiplies local differential
bounds across stages \cite{VahiMirnia2021,VahiMirnia2025}.  These arguments do
not rule out an invariant partition of a complete implemented stage.

The implementation has several such quotients.  These are not low-probability
trails: equality in the quotient is preserved for every input and every key.
A difference-distribution table for individual word differences can therefore
look satisfactory even though a proper partition remains closed.  The
16-bit data path also makes a complete reset codebook cost only \(2^{16}\)
chosen messages.  Finally, the master key is divided into eight disjoint
32-bit pairs, so an outer stage can be removed without guessing any inner key
material.  Similar small-block designs have exhibited related structural
failures \cite{Saarinen2011}.

The main results are as follows.
\begin{enumerate}
\item We derive the implemented diffusion layer and prove that every keyed
stage is triangular.  This gives exact quotient partitions and a direct
counting distinguisher.
\item We characterize the autonomous nibble, enumerate the fibres created by
the lossy key derivation, and show that initialization has 128, rather than
144, reachable bits for each fixed key.
\item We give an inward key-recovery algorithm based on a single pivot and a
proved bidirectional carry filter.  Exhaustive traversal is complete without
a ranking assumption; bounded traversal is measured statistically.
\end{enumerate}

The target throughout is the distributed \code{SEPAR/main.c}.  In particular,
the nibble assignments are sequential and two subkey updates use bitwise OR.
Replacing either operation would define a different cipher.

\section{Cipher and adversarial model}
\label{sec:model}

All additions and subtractions are in \(\W\), while \(\xor\) denotes XOR.
Bits are numbered from zero at the least-significant bit.  Write a word as
nibbles \((a,b,c,d)\), most significant first.  The key is
\(K=(K_1,\ldots,K_8)\), with disjoint, independently used
\(K_i=(k_{i,0},k_{i,1})\in\W^2\).

Let \(S_1,S_2,S_3,S_4\) be the four coordinate S-boxes, and \(\Sigma\) their
parallel application.  In hexadecimal, the implemented tables are
\[
\begin{array}{c|*{16}{c}}
x&0&1&2&3&4&5&6&7&8&9&A&B&C&D&E&F\\\hline
S_1&1&F&B&2&0&3&5&8&6&9&C&7&D&A&E&4\\
S_2&6&A&F&4&E&D&9&2&1&7&C&B&0&3&5&8\\
S_3&C&2&6&1&0&3&5&8&7&9&B&E&A&D&F&4\\
S_4&D&B&2&7&0&3&5&8&6&C&F&1&A&4&9&E
\end{array}
\]
Define
\begin{equation}
G(u)=u\OR\bigl(S_1((u\mathbin{\gg}6)\mathbin{\&}15)\mathbin{\ll}6\bigr).
\label{eq:G}
\end{equation}
At stage \(i\), the source derives
\begin{align}
\kappa_0&=k_{i,0},&\kappa_1&=k_{i,1},\nonumber\\
\kappa_2&=G(\rol(k_{i,0},6))\xor(i+2),&
\kappa_3&=G(\rol(k_{i,1},10))\xor(i+3).
\label{eq:derive}
\end{align}
For \(X_u(x)=x\xor u\), the exact stage is
\begin{align}
\E_i={}&X_{\kappa_2\xor\kappa_3}\circ\Sigma
\circ X_{\kappa_0\xor\kappa_1}\circ L\circ\Sigma
\circ X_{\kappa_3}\circ L\circ\Sigma\nonumber\\
&\circ X_{\kappa_2}\circ L\circ\Sigma
\circ X_{\kappa_1}\circ L\circ\Sigma\circ X_{\kappa_0}.
\label{eq:stage}
\end{align}
Here \(L\) is the diffusion function derived in
Section~\ref{sec:structure}.  All components are bijective;
\code{DEC\_Block} implements \(\D_i=\E_i^{-1}\).

For initialized state words \(s_1,\ldots,s_8\), define
\begin{align}
R_0(x)&=x,\nonumber\\
R_t(x)&=\E_t(R_{t-1}(x)+s_t),\qquad 1\leq t\leq8.
\label{eq:cascade}
\end{align}
The reset one-word ciphertext is \(R_8(x)\).  The in-place state update after
a word matters for multiword transcripts.  The outer experiment requires only
the source identity \(s'_8=s_8+v_{45}\), where \(v_{45}\) is the stage-four
output.

The adversary may choose and reuse a 128-bit IV, reset before each chosen
message, and submit chosen sequences of 16-bit words.  This is a
repeated-reset/chosen-IV model, not a nonce-respecting model; applicability
therefore depends on an interface that permits nonce reuse or arbitrary state
reset.  No decryption oracle is needed.  The outer experiment uses
two-word messages; inward recovery uses complete reset one-word codebooks at
a deterministic public IV family.  Two stateful 64-word messages are reserved
for verification.  The proof-of-concept program contains a local oracle; the
secret is used to generate replies and, when requested, to report experimental
ranks, but it does not enter a search or acceptance decision.

\section{Structural weaknesses}
\label{sec:structure}

\subsection{Sequential linear layer}

The source extracts \((a,b,c,d)\), executes sequentially
\[
a\mathrel{\xor}=c,\quad b\mathrel{\xor}=d,\quad
c\mathrel{\xor}=b,\quad d\mathrel{\xor}=a,
\]
then XORs the packed result with rotations by 12 and 8 bits.

\begin{lemma}[Implemented diffusion]
\label{lem:linear}
\begin{align}
L(a,b,c,d)&=(b\xor c,b,a,a\xor d),\label{eq:L}\\
L^{-1}(A,B,C,D)&=(C,B,A\xor B,C\xor D).
\end{align}
The nibble branch number is two.
\end{lemma}
\begin{proof}
After the four assignments the tuple is
\((a\xor c,b\xor d,c\xor b\xor d,d\xor a\xor c)\).  Substitution into
the two rotations and cancellation gives \eqref{eq:L}; substitution gives the
inverse.  A difference confined to \(d\) remains confined to the fourth
output nibble, attaining branch number \(1+1\).  Bijectivity maps every
nonzero input difference to a nonzero output difference, so no sum below two
is possible.
\end{proof}

One keyed S-box-plus-linear round has dependencies
\begin{align}
a'&=S_2(b\xor\beta)\xor S_3(c\xor\gamma),&
b'&=S_2(b\xor\beta),\nonumber\\
c'&=S_1(a\xor\alpha),&
d'&=S_1(a\xor\alpha)\xor S_4(d\xor\delta).
\label{eq:dependencies}
\end{align}

\begin{theorem}[Triangular stage]
\label{thm:triangular}
For every stage and key there are permutations on \(\Zb4\) such that
\begin{equation}
\E_i(a,b,c,d)=\bigl(P_b(a),Q(b),R_b(c),T_{a,b,c}(d)\bigr).
\label{eq:triangular}
\end{equation}
\end{theorem}
\begin{proof}
Four applications of \eqref{eq:dependencies} yield
\[
\begin{array}{c|cccc}
 &a&b&c&d\\ \hline
1&\{b,c\}&\{b\}&\{a\}&\{a,d\}\\
2&\{a,b\}&\{b\}&\{b,c\}&\{a,b,c,d\}\\
3&\{b,c\}&\{b\}&\{a,b\}&\{a,b,c,d\}\\
4&\{a,b\}&\{b\}&\{b,c\}&\{a,b,c,d\}.
\end{array}
\]
Final S-boxes and XORs add no dependencies.  Bijectivity of \(\E_i\), after
fixing preceding triangular coordinates, forces each one-coordinate component
to be injective and hence a permutation.
\end{proof}

Define
\[
\hi(x)=x\mathbin{\gg}8,\quad
\midb(x)=(x\mathbin{\gg}4)\mathbin{\&}255,\quad
\lane(x)=(x\mathbin{\gg}8)\mathbin{\&}15,\quad
\topx(x)=x\mathbin{\gg}4.
\]

\begin{corollary}[Exact quotient partitions]
\label{cor:quotients}
For each \(\pi\in\{\hi,\midb,\lane,\topx\}\) there is a permutation
\(q_{i,K,\pi}\) satisfying
\begin{equation}
\pi(\E_i(x))=q_{i,K,\pi}(\pi(x))\quad\text{for every }x\in\W.
\label{eq:quotient}
\end{equation}
The same holds for \(\D_i\).  Equality in each quotient is therefore
preserved with probability one in both directions.
\end{corollary}
\begin{proof}
The triangular theorem gives well-defined maps on the partitions.  A quotient
collision would send two equal-size input fibres into one output fibre,
contradicting bijectivity.  Inversion gives the statement for \(\D_i\).
\end{proof}

Thus, for example, an input difference confined to the low byte remains in the
low byte with probability one.  This does not contradict small probabilities
for particular full-word differences; it shows that full-word trails are the
wrong projection.

\subsection{Counting the defect}

There are \(16!\) choices for \(Q\), \((16!)^{16}\) each for \(P_b\) and
\(R_b\), and \((16!)^{16^3}\) for \(T_{a,b,c}\).
\begin{proposition}[Permutation-family bound]
\label{prop:count}
Exactly
\begin{equation}
|\mathcal T|=(16!)^{1+16+16+4096}=(16!)^{4129}
\end{equation}
permutations of \(\W\) have form \eqref{eq:triangular}.  Hence a uniform
random permutation lies in \(\mathcal T\) with probability
\[
\frac{(16!)^{4129}}{65536!}<2^{-771328.03}.
\]
Preserving only the top-12-bit partition already has probability below
\(2^{-729538.24}\).
\end{proposition}
\begin{proof}
Independently chosen component permutations give distinct bijective triangular maps
and every such map reveals its components.  Divide by \(65536!\).  For the
last claim, a partition-preserving permutation has exactly
\(4096!(16!)^{4096}\) choices: permute the 4096 fibres and then permute the 16
points independently inside each fibre.
\end{proof}

\subsection{Autonomous lane and lossy key fibres}
\label{sec:fibres}

Let \(\beta_j\) be bits 11--8 of \(\kappa_j\).  Put
\[
p_0=S_2(b\xor\beta_0),\quad p_1=S_2(p_0\xor\beta_1),\quad
p_2=S_2(p_1\xor\beta_2),\quad p_3=S_2(p_2\xor\beta_3).
\]
The second output nibble is exactly
\begin{equation}
q_{\boldsymbol\beta}(b)=
S_2(p_3\xor\beta_0\xor\beta_1)\xor\beta_2\xor\beta_3.
\label{eq:laneformula}
\end{equation}
Stage constants affect only the low nibble.  We call
\(\nu=(\beta_0,\beta_1,\beta_2,\beta_3)\) a \emph{lane tuple} and write
\(q_\nu\) for its induced permutation.

\begin{lemma}[Lane injectivity; computer-assisted]
\label{lem:laneinjective}
All \(2^{16}\) tuples \(\boldsymbol\beta\) induce distinct permutations in
\eqref{eq:laneformula}.
\end{lemma}
\begin{proof}
The supplied checker evaluates all 16 inputs for every tuple, packs each
output list into 64 bits, sorts the 65,536 signatures, and compares adjacent
values.  This exhausts the finite domain.
\end{proof}

The OR in \eqref{eq:G} is non-injective.  For
\(g(x)=x\OR S_1(x)\),
\[
\begin{array}{c|ccccccccc}
y&1&3&4&7&9&\mathrm B&\mathrm D&\mathrm E&\mathrm F\\ \hline
|g^{-1}(y)|&1&1&1&2&1&1&1&3&5.
\end{array}
\]
Fixing a raw word's direct and derived lane nibbles produces 256 buckets:
128 have size 128, 64 size 192, and 64 size 576.  The two raw words range over
a Cartesian product.

\begin{proposition}[Raw lane fibres]
\label{prop:fibres}
Every lane tuple is realizable.  Its raw 32-bit key fibre has the exact size
distribution in Table~\ref{tab:fibres}; distinct fibres partition all
\(2^{32}\) stage keys.
\end{proposition}
\begin{table}[ht]
\centering
\caption{Exact raw \(K_i\) fibre distribution.}
\label{tab:fibres}
\begin{tabular}{rr}
\toprule
Pairs per lane & Number of lanes \\ \midrule
16,384&16,384\\24,576&16,384\\36,864&4,096\\
73,728&16,384\\110,592&8,192\\331,776&4,096\\ \bottomrule
\end{tabular}
\end{table}
\begin{proof}
For the first raw word, one nibble fixes bits 11--8.  The derived nibble fixes
bits 5--4 and two bits of \(g\) applied to the low nibble; six bits remain
free.  The four two-bit classes contain 2, 3, 2, and 9 inputs, yielding bucket
sizes 128, 192, 128, and 576.  The second word has the same histogram at
rotated positions.  Cartesian products give the table.  Every raw pair maps
to one tuple, proving disjoint coverage.
\end{proof}

\section{Initialization and reachable state}
\label{sec:init}

For one initialization round, compute
\begin{align}
v_1&=\E_1(s_1+s_3+s_5+s_7),\nonumber\\
v_i&=\E_i(v_{i-1}+s_i),\qquad2\leq i\leq8,
\end{align}
then output
\begin{equation}
t_1=s_1+v_8,\qquad t_i=s_i+v_{i-1}\quad(2\leq i\leq8).
\label{eq:initround}
\end{equation}

\begin{theorem}[Explicit initialization inverse]
\label{thm:init}
One round is a permutation of the eight-word state.  Its inverse is
\begin{align}
s_1&=t_1-\E_8(t_8),\nonumber\\
s_i&=t_i-\E_{i-1}(t_{i-1}),\qquad3\leq i\leq8,\nonumber\\
s_2&=t_2-\E_1(s_1+s_3+s_5+s_7).
\label{eq:initinverse}
\end{align}
Thus four rounds biject the 128-bit IV onto the eight final state words.
The initialized LFSR is the deterministic value
\begin{equation}
\ell=\E_8(t_8)\OR\code{0x0100}.
\label{eq:lfsr}
\end{equation}
\end{theorem}
\begin{proof}
For \(2\leq i\leq8\), \eqref{eq:initround} gives
\(v_i=\E_i(t_i)\).  Hence \(s_1=t_1-v_8\) and
\(s_i=t_i-v_{i-1}\) for \(i\geq3\).  These words determine \(v_1\), then
\(s_2=t_2-v_1\).  Substitution proves a two-sided inverse.  In the last round
the value assigned to the LFSR is \(v_8=\E_8(t_8)\), followed by the OR.
\end{proof}

The ninth state word adds no entropy.  Exactly \(2^{128}\) of the
nominal \(2^{144}\) encodings are reachable.  This inverse is also used by
the counterexample certificate below to turn a chosen reachable state into
its unique IV.

\section{Attack overview}
\label{sec:overview}

The attack removes the eight keyed stages in reverse order.  Its two phases
use the same structural defect at different resolutions.

First, fix an IV and encrypt every possible first word followed by seven
short probes.  The first word runs through every possible translation before
the last stage.  The autonomous nibble of that stage therefore gives a small
conjugacy problem: the \(2^{16}\) lane permutations induced by subkey-nibble
tuples can be scored directly from the transcript.  Each such tuple identifies
a disjoint fibre of raw 32-bit stage keys.  Candidates within that fibre are
then ordered at byte and word resolution.  This produces an ordering of
candidates for \(K_8\); it does not assume that the true one is ranked first.

For the remaining stages, collect the complete one-word reset codebook under
several fixed IVs.  Regard each codebook as a table \(C_8\).  To remove stage
\(t\), try a stage key and two state bytes, apply the inverse stage to every
table entry, subtract those bytes, and obtain a candidate table \(C_{t-1}\).
A carry theorem gives four necessary conditions on this table and its inverse.
Only candidates that violate one of these conditions are discarded.  The
same step is repeated until stage one, where a constant-difference identity
tests the final key exactly.  Candidate keys that reach a leaf are checked by
replaying all collected codebooks and two held-out stateful transcripts.

The practical search tries only a few leading score tiers.  The completeness
argument is separate: visit every lane, every raw key in its fibre, and every
state-byte candidate that passes the necessary conditions.  Since scores are
then used solely to choose an order, they cannot affect coverage.  This
distinction is important: Section~\ref{sec:outer} gives a reachable instance
for which the true last-stage lane is not the highest-scoring one.

\section{Recovering the last stage}
\label{sec:outer}

Fix one IV.  For every prefix \(p\in\W\), query the second ciphertext word
\(C_p(q)\) of \((p,q)\) for
\begin{equation}
q\in\{0,1,2,4,8,15,16\}.
\label{eq:probes}
\end{equation}
During encryption of the first word, put
\(v_{12}=\E_1(p+s_1)\) and
\(v_{r,r+1}=\E_r(v_{r-1,r}+s_r)\) for \(2\leq r\leq7\).
The map \(p\mapsto v_{45}\) is a composition of translations and
permutations, hence a bijection.  The source then updates
\(s'_8=s_8+v_{45}\).  Exhausting \(p\) therefore visits every possible outer
translation once, although the other state words remain correlated.

For the second word \(q\), let \(v_{78}(p,q)\) denote its stage-seven output
from the state reached after prefix \(p\), and put
\(z(p,q)=\lane(v_{78}(p,q)+s'_8(p))\).  If \(\nu\) is the true outer lane
tuple, the observed nibble is
\(y(p,q)=\lane(C_p(q))=q_\nu(z(p,q))\).  For
\(\Delta=\{1,2,4,8,15,16\}\), define
\begin{align}
M[a,b]&=\#\{(p,d):y(p,0)=a,\ y(p,d)=b,\ d\in\Delta\},\nonumber\\
H[a,b]&=\#\{(p,d):z(p,0)=a,\ z(p,d)=b,\ d\in\Delta\}.
\end{align}
Then exactly
\begin{equation}
M[q_\nu(a),q_\nu(b)]=H[a,b].
\label{eq:conjugacy}
\end{equation}
The score of candidate lane tuple \(\mu\) is
\begin{equation}
\mathsf{sc}_4(\mu)=\sum_{j=0}^{15}M[q_\mu(j),q_\mu(j+1\bmod16)].
\label{eq:outerS4}
\end{equation}
Writing \(\sigma=q_\nu^{-1}\circ q_\mu\) reduces it to
\begin{equation}
\mathsf{sc}_4(\mu)=\sum_j H[\sigma(j),\sigma(j+1)].
\label{eq:scoreidentity}
\end{equation}
This identity explains the statistic, but it does not make the identity cycle
maximal in \(H\).

The same transcript is also counted at byte and word resolution.  Let
\(M_8[u,v]\) and \(M_{16}[u,v]\) count the observed high-byte and full-word
edges, respectively.  For a raw pair \(K\), put
\(q_K(u)=\hi(\E_{8,K}(2^8u))\); the quotient theorem makes the low byte of
this representative irrelevant.  The two lifting scores are
\begin{align}
\mathsf{sc}_8(K)&=\sum_{u=0}^{255}M_8[q_K(u),q_K(u+1)],\label{eq:outerS8}\\
\mathsf{sc}_{16}(K)&=\sum_{x\in\W}M_{16}[\E_{8,K}(x),\E_{8,K}(x+1)],
\label{eq:outerS16}
\end{align}
where indices wrap around their domains.  These scores refine the order inside
a lane fibre; they do not remove a raw key from exhaustive traversal.

\subsection{Why the best-scoring lane need not be correct}

The following certificate was found by search and then verified
deterministically.  Inverting four initialization rounds reconstructs the IV;
direct computation verifies every entry of \eqref{eq:conjugacy} and enumerates
all \(2^{16}\) lanes.  The key and IV are
\begin{align*}
K={}&\code{E327EAACDFC5765633A8CECCCBC7D7E0}\nonumber\\[-1mm]
&\code{CD7A3ACFFA675B06E4E7004901112801},\\
\mathrm{IV}={}&\code{8D488D31DABD103BC78AB4351BCA28F7}.
\end{align*}
The winning tuple \code{6789} scores 27,899, whereas the true tuple
\code{6781} scores 27,747 and has rank two.  Thus universal rank-one or
maximality claims are false, even on a reachable initialized state.

\subsection{Complete enumeration of \(K_8\)}

We therefore sort all
\(2^{16}\) lane tuples by \eqref{eq:outerS4} with a deterministic tie order.
For rank \(r=1,\ldots,2^{16}\), enumerate the complete preimage fibre of the
lane-to-key map and emit every raw pair in that disjoint fibre.  The high-byte
score orders the fibre; the full-word score refines only its maximal
high-byte tier.  Every unrefined and non-maximal remainder is retained for
later traversal.  Proposition~\ref{prop:fibres} then proves coverage of all
\(2^{32}\) outer pairs.  The rank-two certificate above is also a direct
regression: the first lane omits the true pair, while the second emits
\(K_8=(\code{0111},\code{2801})\).

\section{Peeling the remaining stages}
\label{sec:inward}

Use \(c\) deterministic reset contexts indexed by \(j\), all sharing the key
but having state words \(s_t^{(j)}\).  Write
\[
s_t^{(j)}=b_t^{(j)}+2^8h_t^{(j)},\qquad0\leq b_t^{(j)},h_t^{(j)}<256,
\]
and let \(R_t^{(j)}\) be \eqref{eq:cascade} in that context.

\begin{definition}[Pivot]
At depth \(t\), the search table is
\begin{equation}
C_t^{(j)}(x)=R_t^{(j)}(x)+2^8h_{t+1}^{(j)},\qquad h_9^{(j)}=0.
\label{eq:pivot}
\end{equation}
The root \(C_8^{(j)}=R_8^{(j)}\) is the observed reset codebook.
\end{definition}

Given a candidate pair \(K\) and high byte \(h\), compute
\begin{equation}
W_{t-1}^{(j)}(x)=\D_{t,K}(C_t^{(j)}(x)-2^8h),
\label{eq:inversecandidate}
\end{equation}
then try a low byte \(b\) and recurse on
\begin{equation}
\widetilde C_{t-1}^{(j)}(x)=W_{t-1}^{(j)}(x)-b.
\label{eq:nextpivot}
\end{equation}

\begin{lemma}[Pivot induction]
\label{lem:pivot}
For the true \(K_t,h_{t+1}^{(j)},b_t^{(j)}\),
\(\widetilde C_{t-1}^{(j)}=C_{t-1}^{(j)}\).
\end{lemma}
\begin{proof}
Substitute \eqref{eq:pivot} and
\(R_t=\E_t(R_{t-1}+s_t)\) into \eqref{eq:inversecandidate}:
\[
W_{t-1}=R_{t-1}+b_t+2^8h_t.
\]
Subtracting \(b_t\) gives \eqref{eq:pivot} at depth \(t-1\).
\end{proof}

The two guessed bytes come from consecutive state words: \(h\) is the high
byte of \(s_{t+1}\), already folded into the current pivot, whereas \(b\) is
the low byte of \(s_t\), removed after inversion.  The high byte \(h_t\) that
remains is carried automatically in the next pivot.  This is why one invariant
suffices at every depth.

At \(t=8\), \(h_9=0\) is fixed and only \(b_8\) is tried.  At depths
7 through 2, the same pair and context-specific state bytes are tried across all
contexts before descending to the next key.

\section{Necessary conditions from carries}
\label{sec:carry}

For a permutation \(F\) and quotient \(\pi\in\{\hi,\midb\}\), define
\begin{equation}
\rho_\pi(F)=\max_u
\bigl|\{\pi(F(x)):\pi(x)=u\}\bigr|.
\label{eq:rho}
\end{equation}
Thus \(\rho_\pi(F)\) is the largest number of output quotient classes reached
from any one input class.

\begin{lemma}[One-translation lemma]
\label{lem:translation}
For every modular translation \(\tau_s(x)=x+s\),
\(\rho_{\hi}(\tau_s),\rho_{\midb}(\tau_s)\leq2\), and the same holds for
\(\tau_s^{-1}\).
\end{lemma}
\begin{proof}
Fixing the high byte leaves only the low byte variable.  Addition of the low
byte of \(s\) contributes either zero or one carry to the high byte.  For the
middle byte, write \(x=2^{12}a+16m+d\).  The low nibble contributes at most
one carry into \(m\); the top nibble affects only discarded bits modulo 256.
Negating \(s\) proves the inverse statement.
\end{proof}

\begin{theorem}[Bidirectional retention bound]
\label{thm:carry}
Let \(R_r\) be a cascade of \(r\) translated SEPAR stages and let
\(C_r=R_r+2^8h\).  Then
\begin{equation}
\rho_{\hi}(C_r),\rho_{\midb}(C_r),
\rho_{\hi}(C_r^{-1}),\rho_{\midb}(C_r^{-1})
\leq\min(2^r,256).
\label{eq:fourbounds}
\end{equation}
\end{theorem}
\begin{proof}
Each \(\E_i\) and \(\D_i\) maps a quotient fibre into exactly one quotient
fibre by Corollary~\ref{cor:quotients}; it only relabels fibres.  Each modular
translation can split one fibre into at most two by
Lemma~\ref{lem:translation}.  Induction through \(r\) stages gives \(2^r\),
capped by 256 possible quotient values.  Inverting the cascade reverses the
order and replaces translations by inverse translations, so the same
induction applies.  The final multiple-of-256 translation is itself
quotient-preserving for both listed quotients.
\end{proof}

After a stage is peeled, retain a joint \((h,b)\) candidate only when all four
quantities in \eqref{eq:fourbounds} meet the bound for \(r=t-1\).  Theorem
\ref{thm:carry} guarantees that the true candidate survives.  Wrong
candidates may survive as well: these are necessary conditions, whereas the
scores below only choose an order among survivors.

\subsection{Efficient state-byte enumeration}

\begin{proposition}[Shift invariance of the implemented support statistics]
\label{prop:shiftinvariant}
Fix a pivot table \(C\), a candidate stage permutation \(\E_K\), and write
\[
W_h=\D_K\circ\tau_{-2^8h}\circ C,
\qquad 0\leq h<256.
\]
For each \(\pi\in\{\hi,\midb\}\), both
\(\rho_\pi(W_h)\) and \(\rho_\pi(W_h^{-1})\) are independent of \(h\),
giving four invariant maxima in total.
More strongly, the forward support size of every fixed input row is unchanged,
and the multiset of inverse-row support sizes is unchanged.
\end{proposition}
\begin{proof}
For two guesses \(h,h'\), direct cancellation gives
\[
W_h=A\circ W_{h'},\qquad
A=\D_K\circ\tau_{2^8(h'-h)}\circ\E_K.
\]
Each factor of \(A\) preserves both quotient partitions, so \(A\) induces a
permutation on their fibres.  Postcomposition by \(A\) merely relabels the
output fibres in every forward row and cannot change a row's support size.
Furthermore \(W_h^{-1}=W_{h'}^{-1}\circ A^{-1}\); here \(A^{-1}\) permutes
the input rows, so the inverse-row support multiset is unchanged.  Taking
maxima proves the claim.
\end{proof}

\begin{proposition}[Exact all-byte sweep]
\label{prop:bytesweep}
Let \(F\) be a permutation of \(\W\) and \(F_b=\tau_{-b}\circ F\),
\(0\leq b<256\).  For every \(b\), all four values
\[
 \rho_{\hi}(F_b),\quad \rho_{\midb}(F_b),\quad
 \rho_{\hi}(F_b^{-1}),\quad \rho_{\midb}(F_b^{-1})
\]
and the ordering statistic
\begin{equation}
 U(F_b)=\sum_{u=0}^{255}
 \bigl|\{\hi(F_b(x)):\hi(x)=u\}\bigr|
 \label{eq:statesupport}
\end{equation}
can be computed exactly in \(O(2^{16})\) time and \(O(2^{16})\) memory
for all 256 values of \(b\) together.
\end{proposition}
\begin{proof}
Fix an input high-byte row and, for each output high byte \(q\), record the
minimum and maximum low residues occurring in that row.  After subtracting
\(b\), output row \(q\) occurs precisely when an old-\(q\) residue is at least
\(b\), or an old-\((q+1)\) residue is less than \(b\).  These are two intervals
of \(b\)-values.  A difference-array sweep over the 256 thresholds therefore
gives every forward row support, its maximum, and its contribution to \(U\).
For the middle-byte quotient the identical argument uses the low-nibble
residue; the high nibble of \(b\) only relabels output fibres.

Build \(F^{-1}\) once.  Since \(F_b^{-1}=F^{-1}\circ\tau_b\), increasing \(b\)
by one moves exactly one boundary column between adjacent high-byte input
rows.  Maintaining output-fibre multiplicities gives all inverse high-byte
supports.  The middle-byte calculation moves one boundary column for each
of the 16 top-nibble values and depends only on \(b\bmod16\).  Every update is
an exact insertion or deletion.  The threshold and sliding sweeps, including construction of \(F^{-1}\),
use \(O(2^{16})\) exact updates and \(O(2^{16})\) memory; the stored inverse
dominates the memory bound.
\end{proof}

Before \(b\) is fixed, any threshold on the four maxima treats all 256 values
of \(h\) identically; the forward-row support sum \(U\) is invariant as well.
The exact filter must therefore be evaluated jointly in \((h,b)\).  The
implementation forms \(\tau_{-b}\circ W_h\), checks all four inequalities, and
only then orders the survivors.  This conclusion concerns these support
statistics, not every possible labelled statistic of \(W_h\).

\begin{corollary}[Zero-low-byte orbit]
\label{cor:zeroorbit}
For fixed \(C,t,K\), the 256 maps \(W_h\) from
Proposition~\ref{prop:shiftinvariant} have identical four retention maxima and
identical values of \(U\).  Hence they are either all
inadmissible or all admissible and tied.  In particular, for the true pair
when the true low byte is zero, all 256 high-byte guesses survive and tie with
the true guess.
\end{corollary}
\begin{proof}
For \(b=0\), the joint candidate after low-byte subtraction is \(W_h\).
Proposition~\ref{prop:shiftinvariant} makes the four maxima independent of
\(h\).  It also preserves the support
size of each forward high-byte row, so their sum \(U\) is independent of
\(h\).  For the true pair and low byte, Theorem~\ref{thm:carry} makes the true
high byte admissible, and therefore makes the whole orbit admissible.
\end{proof}

\section{Candidate ordering}
\label{sec:ranking}

At each depth the search must order three finite sets: autonomous nibble
permutations, raw key pairs within a nibble fibre, and surviving state-byte
guesses.  The first two scores count adjacent transitions in a quotient of the
observed codebook.  The last prefers candidate tables with smaller high-byte
support.  None of these scores rejects a candidate.
Candidates with the same score form a \emph{score tier}; a tier budget always
visits the whole tie at its boundary.

Let \(J\) be the current set of active recovery contexts.  For a pivot table
\(C\), a quotient \(\pi\), and the fixed difference set \(\Delta\), form the
edge matrix
\begin{equation}
M_{\pi,C}[u,v]=\#\{(x,d):d\in\Delta,\ \pi(C(x))=u,
 \pi(C(x+d))=v\}.
\label{eq:edges}
\end{equation}
For a lane tuple \(\nu\), lane offset \(\lambda\), and raw pair
\(K\), let
\begin{align}
\mathsf{sc}_4(C;\nu,\lambda)
 &=\sum_{z=0}^{15}
 M_{\lane,C}[q_\nu(z)+\lambda,q_\nu(z+1)+\lambda],
 \label{eq:inS4}\\
Q_K(z)&=\hi(\E_K(2^8z)),\nonumber\\
\mathsf{sc}_8(C;K,h)
 &=\sum_{z=0}^{255}
 M_{\hi,C}[Q_K(z)+h,Q_K(z+1)+h].
\label{eq:inS8}
\end{align}
Here \(2^8z\) is simply a representative of the high-byte fibre \(z\);
Corollary~\ref{cor:quotients} makes the choice of its low byte immaterial.
Indices are reduced modulo 16 or 256.  The implemented lane and pair scores
are exactly
\begin{align}
\mathsf{Sc}_4(\nu)&=\sum_{j\in J}\max_{0\leq\lambda<16}
 \mathsf{sc}_4(C_t^{(j)};\nu,\lambda),\label{eq:aggregateS4}\\
\mathsf{Sc}_8(K)&=\sum_{j\in J}\max_{0\leq h<256}
 \mathsf{sc}_8(C_t^{(j)};K,h).\label{eq:aggregateS8}
\end{align}
Thus every lane score examines all 16 translations and every pair score
examines all 256 high-byte translations.  There is no translation budget.

For a pair \(K\), define
\[
H_{K,j}=\operatorname*{arg\,max}_{0\leq h<256}
 \mathsf{sc}_8(C_t^{(j)};K,h).
\]
Let
\[
M_{16,C}[u,v]=\#\{(x,d):d\in\Delta,\ C(x)=u,\ C(x+d)=v\},
\]
and, with \(x+1\) reduced modulo \(2^{16}\), define
\begin{equation}
\mathsf{sc}_{16}(C;K,h)=\sum_{x\in\W}
 M_{16,C}[\E_K(x)+2^8h,\E_K(x+1)+2^8h].
\label{eq:inS16}
\end{equation}
The implementation computes this sparse full-word score only for
\(h\in H_{K,j}\), and uses the aggregate
\begin{equation}
\mathsf{Sc}_{16}(K)=\sum_{j\in J}\max_{h\in H_{K,j}}
\mathsf{sc}_{16}(C_t^{(j)};K,h).
\label{eq:aggregateS16}
\end{equation}
This \(\mathsf{Sc}_{16}\) value orders candidates
\emph{inside one complete tie of \(\mathsf{Sc}_8\)}.  It neither defines a
fourth budget nor removes a candidate: every member of a visited
\(\mathsf{Sc}_8\)-tier is tried, regardless of its \(\mathsf{Sc}_{16}\) order.

For a fixed pair \(K\) and context \(j\), define
\begin{equation}
 T_{t,K}^{(j)}(h,b)=
 \tau_{-b}\circ\D_{t,K}\circ\tau_{-2^8h}\circ C_t^{(j)}
 \label{eq:jointcandidate}
\end{equation}
and its exact admissible set
\begin{equation}
 \mathcal A_{t,K}^{(j)}=
 \{(h,b):\rho_\pi(T_{t,K}^{(j)}(h,b)),
 \rho_\pi((T_{t,K}^{(j)}(h,b))^{-1})
 \leq\min(2^{t-1},256)\ \text{for }\pi\in\{\hi,\midb\}\}.
 \label{eq:admissible}
\end{equation}
Here \(h=h_{t+1}^{(j)}\) and \(b=b_t^{(j)}\).  The packed value
\(2^8h+b\) only encodes a candidate pair of bytes; it is not one SEPAR state
word.  Members of \(\mathcal A\) are ordered by increasing
\(U(T_{t,K}^{(j)}(h,b))\) from \eqref{eq:statesupport}.  At the root, \(h_9=0\)
is fixed, so the analogous list contains only admissible \(b_8\).

The inward search has three tier boundaries; the outer phase has a separate
limit on last-stage lanes:
\begin{enumerate}
\item visit score tiers of the \(2^{16}\) lanes under
\(\mathsf{Sc}_4\);
\item expand each visited lane into its exact raw-pair fibre and visit score
tiers under \(\mathsf{Sc}_8\), with \(\mathsf{sc}_{16}\) used only as a within-tier order;
\item for each pair and active context, visit score tiers of its exact
\((h,b)\) survivors under \(U\), or of \(b_8\) at the root.
\end{enumerate}
The three budgets count distinct score values, rather than individual
candidates.  A zero budget denotes all tiers.  With all three budgets zero,
the candidate set is independent of every score and tie order.

\subsection{Context retirement in bounded search}
\label{sec:retirement}

Corollary~\ref{cor:zeroorbit} explains why expanding every tied state
candidate can dominate a bounded run.  The implementation therefore applies
one deterministic rule when the state-tier budget is nonzero.
When the lowest-\(U\) tier for the current context contains more than one
candidate and at least two contexts remain active, that context is retired for
the remainder of the current branch.  The root applies the same rule to its
lowest-support \(b_8\) tier.  A retired context contributes neither later
ranking scores nor reconstructed state constraints; its status is restored on
backtracking.  At least one context always remains active.  All input to this
decision is computed from public codebooks and the current public candidate.

\begin{lemma}[Exact-filter monotonicity under retirement]
\label{lem:retirement}
Let \(J^{\prime}\subseteq J\).  On the branch containing the true key and true
pivots, restricting the componentwise exact tests from \(J\) to
\(J^{\prime}\) cannot reject the true key path on any retained context.  In
particular, retirement cannot invalidate pivot induction or the carry and
stage-one tests that remain.
\end{lemma}
\begin{proof}
For every \(j\in J\), Lemma~\ref{lem:pivot} supplies the true next pivot and
Theorem~\ref{thm:carry} places its true \((h,b)\) in
\(\mathcal A_{t,K}^{(j)}\).  The exact multi-context condition is a conjunction,
equivalently a Cartesian product of the per-context survivor sets.  Deleting a
component leaves the restriction of the true vector in the remaining product.
The same argument applies at the root, and the true stage-one factor satisfies
\eqref{eq:factor} separately in every retained context.
\end{proof}

Lemma~\ref{lem:retirement} is not a theorem about bounded ranking.  After a
retirement, the program rebuilds edge tables and recomputes
\(\mathsf{Sc}_4\) and \(\mathsf{Sc}_8\) over the smaller active set.  A true lane,
pair, or state coordinate may then lie beyond a nonzero tier budget.  Moreover,
the rule can retire a context without knowing the rank of its true coordinate.
The success probability and runtime of this bounded, dynamically re-ranked
algorithm must therefore be evaluated end to end.

In exhaustive search the state budget is zero, so retirement is disabled.
Every context remains active, every score tier is visited, and changing a
score or tie order changes only the enumeration order, not the finite set of
leaves.

\subsection{Stage-one factor}

At the final pivot,
\begin{equation}
C_1^{(j)}(x)=\E_{1,K_1}(x+s_1^{(j)})+2^8h_2^{(j)}.
\end{equation}
For every candidate \((K,h)\), test whether
\begin{equation}
\D_{1,K}(C_1^{(j)}(x)-2^8h)-x
\label{eq:factor}
\end{equation}
is constant for all \(2^{16}\) inputs.  If so, the constant is the candidate
\(s_1^{(j)}\).  The implementation enumerates all 256 values of \(h\) in
every context, rather than retaining the first factor.  For the true pair and
true \(h_2^{(j)}\), \eqref{eq:factor} is identically \(s_1^{(j)}\).

\section{Completeness}
\label{sec:correctness}

\begin{theorem}[Eventual completeness and order independence]
\label{thm:complete}
For every key, deterministic IV family, and transcript produced by the
implemented cipher, consider the traversal that enumerates all last-stage
lanes and their raw-key fibres, uses zero inward tier budgets, never retires a
context, and continues after accepted leaves until every finite loop is
exhausted.  It visits the original numeric key and its initialized states
after finitely many candidates, and that leaf reproduces the transcript.
Replacing an ordering score or a tie order changes only the order of this
traversal, not its coverage.
\end{theorem}
\begin{proof}
The outer traversal visits all \(2^{16}\) lane tuples and every raw pair in
each disjoint fibre, so Proposition~\ref{prop:fibres} guarantees that it
supplies the true \(K_8\).  All three inward tier budgets are zero and
retirement is disabled.  For the true \(K_8\), Lemma~\ref{lem:pivot}
gives \(\D_{8,K_8}(C_8^{(j)})-b_8^{(j)}=C_7^{(j)}\), and
Theorem~\ref{thm:carry} with \(r=7\) makes every true \(b_8^{(j)}\) admissible.
The root visits every admissible low byte with fixed \(h_9^{(j)}=0\), and hence
reaches all true \(C_7^{(j)}\) pivots.  At each stage \(t=7,\ldots,2\), all lane
tuples and every pair in their fibres are visited.  For the true pair,
Lemma~\ref{lem:pivot} maps the true bytes
\((h_{t+1}^{(j)},b_t^{(j)})\) to the true next pivot, while
Theorem~\ref{thm:carry} puts that byte pair in the exactly computed
\(\mathcal A_{t,K}^{(j)}\).  A zero state-tier budget visits every member of
every such set.  Induction therefore reaches the true stage-one pivots in all
contexts.  Complete lane and pair traversal visits the true \(K_1\), all 256
values of every \(h_2^{(j)}\) are tested, and \eqref{eq:factor} returns each
true \(s_1^{(j)}\).

Every loop ranges over a finite set.  Scores are used only to sort or partition
sets that the exhaustive search traverses in full, so altering their order does not
alter coverage.  The true key and states reproduce initialization, every
cached inward codebook, and both held-out transcripts, hence the original-key
leaf is accepted.
\end{proof}

The proof concerns a traversal that continues after finding a candidate.  The
supplied program instead stops at the first key that reconstructs the active
initializations and reproduces every collected codebook and both held-out
64-word transcripts.  Such a key agrees with those finite transcripts.
The inward leaf check does not replay the outer prefix transcript; that
transcript orders candidates and supplies \(K_8\).
The theorem guarantees that continued exhaustive enumeration reaches the
original key; it does not assert that a finite transcript uniquely identifies
its numerical representation.

This is an eventual enumeration theorem, not a polynomial-time or
minute-scale bound.  For \(c\) contexts, a deliberately coarse unfiltered
leaf bound is
\begin{equation}
2^{32}\,2^{8c}\bigl(2^{32+16c}\bigr)^6 2^{32+8c}
=2^{256+112c}.
\label{eq:treecoarse}
\end{equation}
The factors count the outer pair, root low bytes, six intermediate pairs with
two state bytes per context, and the stage-one pair and high bytes.  The bound
proves finiteness but is too loose to guide an implementation; exhaustive search over the
\(2^{256}\) master keys is a tighter unconditional fallback.  The practical
claim concerns bounded score tiers and is therefore statistical.

\section{Data and computational complexity}
\label{sec:effort}

The outer experiment contains exactly
\begin{equation}
7\cdot2^{16}=458{,}752
\label{eq:outerqueries}
\end{equation}
separately reset chosen-plaintext messages of two words each, hence 917,504
online word encryptions.  It
contributes \(6\cdot2^{16}=393{,}216\) directed nonzero-difference edges.

For each of \(c\) deterministic IVs, inward recovery resets before every input
\(x\) and collects a \(2^{16}\)-entry one-word codebook.
It therefore uses \(c2^{16}\) reset messages and the same number of word
encryptions.  Verification adds two held-out reset messages, each carrying a
stateful 64-word transcript under a distinct IV.  Every response is cached
before search, so backtracking and failed leaves require no further queries.
For \(c=8\), the stand-alone inward phase uses 524,290 messages containing
524,416 words.  If the two phases are collected separately, their sum is
983,042 messages containing 1,441,920 words.
Choose the fixed outer IV as the first member of the inward IV family.  Each
outer two-word response already contains the first ciphertext word
\(R_8(p)\), so retaining one response for every prefix supplies that inward
codebook at no extra query cost.  A cached implementation therefore uses
\begin{align}
Q_{\rm msg}&=458{,}752+(c-1)2^{16}+2,\nonumber\\
Q_{\rm word}&=917{,}504+(c-1)2^{16}+128.
\label{eq:uniquequeries}
\end{align}
For \(c=8\), these are 917,506 messages containing 1,376,384 words.  All
responses are cached before the search, so neither count has a per-leaf
factor.

A black-box implementation can collect the fixed outer experiment, the
remaining codebooks, and the held-out transcripts once, then cache them for
the whole search.  The dominant offline costs are summarized below.  Here
\(N\) is the number of raw pairs in the selected outer fibres, \(B\) is the
number given a full-word outer score, and \(a=|J|\) is the number of active
inward contexts.

\begin{table}[ht]
\centering
\caption{Dominant offline costs at each level of the search.}
\label{tab:work}
\small
\begin{tabular}{ll}
\toprule
Operation & Work \\ \midrule
Outer high-byte lift & \(2^8N\) stage evaluations \\
Outer full-word refinement & \(2^{16}B\) stage evaluations \\
Inward lane ordering & \(a2^{24}\) matrix additions \\
Inward pair ordering in a fibre of size \(N_\nu\) &
  \(2^8N_\nu\) stage calls and \(a2^{16}N_\nu\) additions \\
Joint state-byte sweep, per pair and context &
  \(2^{24}\) inverse-stage calls and \(O(2^{24})\) updates \\
Root low-byte sweep, per context & \(O(2^{16})\) \\
Stage-one factor test, per \((K,h)\) & at most \(2^{16}\) inverse-stage calls \\
\bottomrule
\end{tabular}
\end{table}

One lane contains at most 331,776 raw pairs.  Enumerating every outer lane
gives coarse bounds \(2^{40}\) and \(2^{48}\) for the high-byte and full-word
passes.  The joint state-byte sweep dominates a typical bounded inward node.
Its 256 high-byte guesses are divided into contiguous ranges processed by
native threads and merged in a deterministic order.  These are costs per
visited candidate or tier; the number of branches reached by bounded search
is an empirical quantity.

\section{Experimental evaluation}
\label{sec:experiments}

The proof of concept is a single portable threaded C program that runs the
last-stage prefix experiment and passes its selected candidate directly to
inward recovery.  Its cipher core is a literal implementation of the reference
operations.  Self-tests
check the published stream, block inversion, the lane formula and exhaustive
lane injectivity, the optimized carry computation against a direct reference,
pivot induction, and the stage-one factor test.  A bounded search that exhausts
its configured tiers reports \code{INCONCLUSIVE}, rather than claiming
failure of the full attack.

The 100-key evaluation below timed the outer and inward phases separately so
that their ranking behaviour could be distinguished.  The artifact's unified
entry point invokes the same two routines in sequence for a complete demo.

We evaluated 100 independently generated 256-bit keys, fixed before the runs.
The inward experiment used eight deterministic IV contexts, seed one, a
600-second timeout, and tier budgets
\[
(\text{lane},\text{pair},\text{state})=(1,1,1).
\]
It was conditioned on the correct externally supplied \(K_8\), so it measures
the inward phase rather than time spent trying earlier last-stage candidates.
A run counted as successful only if it passed transcript verification and
returned the sampled numeric key.  Every timeout, inconclusive run, or other
output counted as a failure.  The outer experiment was then evaluated on the
same keys and fixed IV, retaining four lane ranks.

The inward phase recovered 95 keys, with three inconclusive runs and two
timeouts.  The outer phase placed the correct \(K_8\) first for all 100 keys.
Consequently the matched event ``outer first and inward success'' also
occurred 95 times.  Table~\ref{tab:bounded-results} gives exact
Clopper--Pearson intervals \cite{ClopperPearson1934}.  The reported 95th
percentile is the nearest-rank order statistic.

\begin{table}[ht]
\centering
\caption{Bounded-search results on the same 100 keys.  Inward recovery is
conditional on the correct external \(K_8\); its timing columns describe the
95 successful runs.}
\label{tab:bounded-results}
\small
\begin{tabular}{lrrrr}
\toprule
Experiment & Exact recovery & Exact 95\% interval & Median & \(\widehat q_{.95}\) \\ \midrule
Inward, given \(K_8\) & 95/100 & [0.8872, 0.9836] & 181.1 s & 284.1 s \\
Outer candidate first & 100/100 & [0.9638, 1] & --- & --- \\
Matched joint event & 95/100 & [0.8872, 0.9836] & --- & --- \\
\bottomrule
\end{tabular}
\end{table}

The exact one-sided 95\% lower confidence bound for inward and matched success
is 0.8977; the corresponding upper bound on failure is 0.1023.  Successful
inward runs ranged from 73.5 to 326.9 seconds, with mean 183.4 seconds.  Across
all trials, including the two timeouts at their observed duration, the median
was 181.0 seconds and the nearest-rank 95th percentile was 285.7 seconds.

The machine was an AMD Ryzen~9 8945HS with 16 logical CPUs, running
Linux~6.18.33.2 under WSL2, GCC~11.4.0, and glibc~2.35.  The evaluator ran two
processes concurrently with eight native threads per process; the reported
times therefore include their contention.  They measure conditional inward
recovery and verification, not data acquisition or end-to-end latency.

Under the model that the sampled keys are independent uniform draws and
the resulting per-key strict-success indicators are Bernoulli trials, failure
probability \(p_f\) with \(f\) failures among \(n\) keys has the exact
one-sided Clopper--Pearson bound, for \(f<n\),
\begin{equation}
p_f\leq \operatorname{Beta}^{-1}_{1-\alpha}(f+1,n-f)
\label{eq:cp}
\end{equation}
at confidence \(1-\alpha\) \cite{ClopperPearson1934}.  The success lower bound
is its complement.  Only one indicator per key enters this calculation;
stage, context, and retirement ranks within a run are correlated and can be
censored by early termination.

For all five missed keys, the true inward S-box lane lay outside the one-tier
budget.  It occurred
in tier 3 at stage 6 for trial 15, and in tiers 8, 9, 4, and 2 at stage 7 for
trials 50, 54, 89, and 91.  Whenever the corresponding branch was reached,
every observed true byte, raw pair, and joint state-byte coordinate ranked
first.  These within-run ranks are correlated and censored by early stopping;
they explain the observed failures but are not independent probability
estimates.

The outer result does not establish universal rank one; the certificate in
Section~\ref{sec:outer} disproves such a statement.  The matched 95/100 result
is the Boolean intersection of the two outcomes on identical keys, not a
product of estimated rates.  It also does not measure the time required to run
inward recovery on false \(K_8\) candidates that precede the true pair.

\section{Discussion and conclusion}

SEPAR's sequential assignments collapse the intended diffusion layer and make
every keyed stage triangular.  Four proper coordinate partitions are
preserved exactly, one nibble evolves autonomously, and only a tiny structured
subset of all 16-bit permutations can occur.  The OR-based key derivation and
the use of disjoint key pairs make these properties exploitable.  Separately,
initialization reaches only 128 bits of the nominal 144-bit state.

The resulting attack is a factorization of the cipher from the outside in.
It enumerates last-stage key fibres, maintains one pivot while peeling the
remaining stages, discards only candidates that fail a proved carry bound, and
closes with an exact stage-one identity.  An exhaustive traversal is complete
and independent of its scores.  The bounded version replaces exhaustive
ordering by three small score-tier budgets and context retirement; its
minute-scale behavior is therefore supported by the experiment, not by an
unproved maximality assumption.

The result is limited to the repeated-reset, chosen-IV setting of
Section~\ref{sec:model}.  We do not give a useful universal sub-\(2^{256}\)
time bound, so the completeness theorem does not imply that every key is
recovered in minutes.  The exhaustive proof and the measured score ordering
are therefore separate claims.

A redesign must replace the linear layer with one having a justified branch
number and no closed coordinate partition, state simultaneous operations
unambiguously, avoid the lossy OR derivation, and diffuse key material between
stages.  A 16-bit external permutation also makes complete codebooks cheap, so
enlarging only the internal state is not sufficient.  Statistical output tests
cannot substitute for this structural analysis.

The accompanying artifact contains the unified C program, the reference cipher
source, the complete 100-key result table, and representative execution logs.

\appendix
\section{Attack traversal pseudocode}

Here \textsc{VisitTiers} visits every row in each selected score tier, including
the whole boundary tie; budget zero selects all tiers.  Lane and pair tiers are
visited in decreasing score order, while state tiers are visited in increasing
\(U\) order.  \textsc{StateList} enumerates all \(2^{16}\) joint byte pairs,
discards only failures of the four proved bounds, and sorts the survivors by
increasing \(U\).  \textsc{RootList} does the same for the 256 low bytes with
\(h_9=0\).  Sets \(J\) are immutable in the pseudocode, so return from a call
automatically restores a retired context.
\begin{verbatim}
Recover(t, C[1..c], J):
  LaneScore(nu) = sum(j in J) max(lambda=0..15) score4(C[j],nu,lambda)
  for lane tier LN in VisitTiers(all nu by decreasing LaneScore, lane_budget):
    for nu in every row of LN:
      PairScore(K) = sum(j in J) max(h=0..255) score8(C[j],K,h)
                     for every K in ExactFibre(nu)
      for pair tier PK in VisitTiers(all K by decreasing PairScore, pair_budget):
        order every K in PK by decreasing score16; do not remove any K
        for K in that order:
          if t == 1:
            Factor(1, K, C, J)
          else:
            Advance(t, K, 1, C, J)

Advance(t, K, j, C, J):
  if j > c: Recover(t-1, C, J); return
  if j not in J: Advance(t, K, j+1, C, J); return
  A = StateList(t, K, C[j])
  F = the complete first score tier of A, if A is nonempty
  if A is nonempty and state_budget > 0 and |F| > 1 and |J| > 1:
    Advance(t, K, j+1, C, J - {j}); return
  for state tier T in VisitTiers(A, state_budget):
    for (h,b) in every row of T:
      Cnext[j] = D_t,K(C[j]-256*h)-b
      Advance(t, K, j+1, C with C[j]=Cnext[j], J)

Factor(j, K1, C, J):
  if j > c:
    if VerifyActiveInitialization() and VerifyEveryCachedCodebook()
       and VerifyBothHeldOutTranscripts():
      emit candidate; executable returns, conceptual complete search continues
  else if j not in J:
    Factor(j+1, K1, C, J)
  else:
    for h = 0..255:
      if D_1,K1(C[j](x)-256*h)-x is one constant s1 for every x:
        record s1 and h; Factor(j+1, K1, C, J)

RootAdvance(j, C, J):
  if j > c: Recover(7, C, J); return
  if j not in J: RootAdvance(j+1, C, J); return
  A = RootList(K8, C[j])
  F = the complete first score tier of A, if A is nonempty
  if A is nonempty and state_budget > 0 and |F| > 1 and |J| > 1:
    RootAdvance(j+1, C, J - {j}); return
  for state tier T in VisitTiers(A, state_budget):
    for b in every row of T:
      Cnext[j] = D_8,K8(C[j])-b
      RootAdvance(j+1, C with C[j]=Cnext[j], J)

Inward(K8):
  C8[j] = cached complete reset codebook for every j
  RootAdvance(1, C8, {1,...,c})

FullAttack:
  for nu in all 65536 outer lanes, ordered by the outer lane score:
    for K8 in every raw pair of ExactFibre(nu), ordered by byte/word scores:
      Inward(K8)
\end{verbatim}

With all three budgets zero, retirement and score-tier pruning are absent, but
the executable still stops at its first accepted leaf.  For the coverage
theorem, the accepted-leaf stop is disabled and the outer loop continues after
each inward success.  It therefore streams every raw pair from every lane fibre
to a non-stopping inward traversal.  With nonzero budgets, the same pseudocode
is a dynamically re-ranked heuristic whose success rate must be measured.

\small
\bibliographystyle{plain}
\bibliography{references}
\end{document}
